\part{Booting}

\chapter{From the power button to the first instruction}

\begin{goals}
  \item What the firmware does before anything of yours exists
  \item How the BIOS decides where to boot from
  \item What the \hexa{AA55} signature is and why nothing boots without it
\end{goals}

\section{The first few milliseconds}

You press the button. The processor initialises with fixed values: \reg{CS} holds \hexa{F000},
\reg{EIP} holds \hexa{FFF0}, and it is in \term{real mode}, behaving like an 8086 from 1978.

At the resulting address there is a jump into the BIOS code. From there, the firmware does its
work:

\begin{enumerate}
  \item \term{POST} (Power-On Self Test): check that RAM responds, that a keyboard exists, that
        disks answer.
  \item Initialise the basic controllers: interrupt controller, timer, disk controller.
  \item Install its own service table in low memory, so software can ask it for things.
  \item Look for a device to boot from.
\end{enumerate}

\section{How the boot device is chosen}

The BIOS walks the devices in the configured order. From each one it reads the first sector ---
512 bytes --- and checks exactly one thing: that \hi{the last two bytes are \hexa{55} and
\hexa{AA}} (which read as a little-endian 16-bit number form \hexa{AA55}).

If the signature is there, the BIOS considers that sector bootable: it copies it to address
\hexa{7C00} and jumps there. If not, it moves to the next device.

\begin{concept}{The contract is two bytes long}
That is the entire contract between firmware and your operating system: 512 bytes ending in
\hexa{AA55}. There is no file format, no header, no integrity check. The BIOS loads those bytes and
jumps. What happens next is entirely your responsibility.
\end{concept}

That is why the last line of the lytlnyblOS bootloader is exactly this:

\filetag{lytlnybl/boot/bootloader.asm}
\begin{lstlisting}[style=asm]
times 510-($-$$) db 0   ; pad with zeros up to byte 510
dw 0xAA55               ; the magic signature in bytes 510 and 511
\end{lstlisting}

\cod{\$} means ``the current address'' and \cod{\$\$} ``the start address of the section'', so
\cod{\$-\$\$} is ``how many bytes I have written so far''. The directive pads with zeros up to 510
and then places the signature. If your code went past 510 bytes, NASM would error out on a negative
count: it is a hard limit.

\begin{warn}{Why \hexa{7C00} and nowhere else}
The address \hexa{7C00} is a relic of the original IBM PC 5150, which had 32 KB of RAM. The
designers wanted the boot sector as high as possible without touching the last kilobyte (reserved),
and \hexa{7C00} is exactly $32\,\text{KB} - 1024$. Forty-five years later, that number is still
there.
\end{warn}

% ============================================================
\chapter{The 512 bytes of the bootloader}

\begin{goals}
  \item Read the complete lytlnyblOS bootloader and understand every line
  \item Understand 16-bit segmentation and the \cod{segment$\times$16 + offset} calculation
  \item Know why the bootloader has to care about the \reg{DS} register
\end{goals}

\section{The whole program}

Before dissecting it, look at the overall shape. The lytlnyblOS bootloader does four things, in
this order:

\begin{center}
\begin{tikzpicture}[font=\sffamily\scriptsize,
  b/.style={draw=accent,fill=accentlight,rounded corners=3pt,minimum width=3.2cm,minimum height=0.95cm,align=center}]
  \node[b] (a) {Say hello\\{\tiny print\_string}};
  \node[b,right=0.5cm of a] (c) {Ask for the\\memory map\\{\tiny int 15h / E820}};
  \node[b,right=0.5cm of c] (d) {Load the kernel\\{\tiny int 13h}};
  \node[b,right=0.5cm of d] (e) {Jump to it\\{\tiny jmp 0900h:0000}};
  \foreach \x/\y in {a/c,c/d,d/e} \draw[-{Stealth[length=2.2mm]},accent,thick] (\x)--(\y);
\end{tikzpicture}
\end{center}

\filetag{lytlnybl/boot/bootloader.asm}
\begin{lstlisting}[style=asm]
start:
    mov ax, 07C0h
    mov ds, ax          ; DS = 0x07C0, so that labels resolve correctly

    mov si, title_string
    call print_string

    mov si, message_string
    call print_string

    call store_memory_map
    call load_kernel_from_disk

    jmp 0900h:0000      ; jump to the kernel, loaded at physical address 0x9000
\end{lstlisting}

Twelve lines and the system is already in motion. Let us understand the first two, which are the
most confusing.

\section{16-bit segmentation: why \hexa{07C0}}

In real mode the processor has 16-bit registers. With 16 bits you can only count to 65,535, that
is 64 KB. But the 8086 wanted to address 1 MB. Intel's solution was to combine \emph{two} 16-bit
registers:

\begin{center}
\begin{tikzpicture}[font=\sffamily\small]
  \node[draw=purple,fill=purplelight,rounded corners=2pt,minimum width=2.9cm,minimum height=0.8cm] (seg) at (0,0) {segment};
  \node[font=\Large,text=gray] at (1.95,0) {$\times\,16$};
  \node[font=\Large,text=gray] at (3.4,0) {$+$};
  \node[draw=accent,fill=accentlight,rounded corners=2pt,minimum width=2.6cm,minimum height=0.8cm] (off) at (5.5,0) {offset};
  \node[font=\Large,text=gray] at (7.4,0) {$=$};
  \node[draw=green,fill=greenlight,rounded corners=2pt,minimum width=3.6cm,minimum height=0.8cm] (fis) at (10.0,0) {physical address};
\end{tikzpicture}
\end{center}

Multiplying by 16 is the same as shifting four bits left, that is, adding a hex zero. So:

\begin{center}
\cod{07C0} $\times$ 16 = \cod{7C00} \qquad and \qquad \cod{7C00} + \cod{0000} = \cod{7C00}
\end{center}

Which is exactly where the BIOS left our code. By setting \reg{DS} = \hexa{07C0}, any label such as
\cod{title\_string} (which the assembler turns into a small offset, say \hexa{6E}) resolves to
physical address \hexa{7C6E}: the right place.

\begin{danger}{If you forget to set DS}
If you do not initialise \reg{DS}, it holds whatever the BIOS left there, usually 0. Then
\cod{title\_string} would resolve to physical address \hexa{6E}, which is in the middle of the
interrupt vector table. You would print garbage --- or hang the machine. It is the number one
mistake of anyone writing their first bootloader, and it produces no message at all: it simply does
not work.
\end{danger}

\begin{warn}{Same place, many names}
An odd consequence of this scheme is that one physical address has many representations:
\cod{07C0:0000}, \cod{0000:7C00} and \cod{0700:0C00} all point at \hexa{7C00}. It is not a design
error; it is what happens when segments overlap every 16 bytes.
\end{warn}

\section{Printing without an operating system}

There is no video driver yet, so we use a BIOS service. They are invoked with \cod{int}, after
putting the function number in \reg{AH}.

\filetag{lytlnybl/boot/bootloader.asm}
\begin{lstlisting}[style=asm]
print_string:
    mov ah, 0Eh         ; video service function 0x0E: teletype output
print_char:
    lodsb               ; AL = [DS:SI], then SI = SI + 1
    cmp al, 0           ; end of string?
    je printing_finished
    int 10h             ; BIOS video service: print AL
    jmp print_char
printing_finished:
    mov al, 10d         ; line feed
    int 10h
    mov ah, 03h         ; read cursor position
    mov bh, 0
    int 10h
    mov ah, 02h         ; set cursor position
    mov dl, 0           ; column 0
    int 10h
    ret
\end{lstlisting}

The \cod{lodsb} instruction is an x86 shorthand: load into \reg{AL} the byte pointed at by
\cod{DS:SI} and increment \reg{SI}. It is exactly the equivalent of \cod{c = *p++;} in C, in a
single instruction.

The end of the routine has an amusing historical detail: code 10 is only a \emph{line feed}, which
moves down a row but does not return to the left margin. So afterwards you must read the cursor
position and set the column to 0 by hand. A direct inheritance from typewriters.

\exercise{Why does \cod{print\_string} use \reg{SI} and not, say, \reg{BX} to walk the string?
Hint: look at which register \cod{lodsb} uses implicitly.}

% ============================================================
\chapter{Reading the disk and loading the kernel}

\begin{goals}
  \item What CHS addressing is and why it still exists
  \item How to ask the BIOS to read sectors
  \item Why the kernel is loaded at \hexa{9000}
\end{goals}

\section{Service \hexa{13}}

The BIOS offers disk access through interrupt \hexa{13}. Function 2 reads sectors, and expects its
parameters spread across half the register file:

\filetag{lytlnybl/boot/bootloader.asm}
\begin{lstlisting}[style=asm]
load_kernel_from_disk:
    mov ax, 0900h
    mov es, ax          ; ES = 0x0900 -> destination physical address 0x9000

    mov ah, 02h         ; function 2: read sectors
    mov al, 60          ; how many sectors to read

    mov ch, 0h          ; cylinder 0
    mov cl, 02h         ; sector 2 (counting starts at 1!)
    mov dh, 0h          ; head 0
    mov dl, 80h         ; drive 0x80 = first hard disk

    mov bx, 0h          ; destination offset -> ES:BX = 0900:0000

    int 13h
    jc kernel_load_error  ; the BIOS sets the carry flag on failure
    ret
\end{lstlisting}

\subsection{CHS: cylinder, head, sector}

Old disks were spinning platters and were addressed by their physical geometry: \term{cylinder}
(which ring), \term{head} (which side of the platter) and \term{sector} (which slice of the ring).
Today disks are electronic and that geometry is fictional, but the interface persists for
compatibility.

\begin{warn}{Sectors start at 1}
In CHS, cylinders and heads are numbered from 0, but \hi{sectors are numbered from 1}. So
\cod{cl = 2} means ``the second sector of the disk'', the first one \emph{after} the boot sector.
This inconsistency has caused countless bugs since 1983.
\end{warn}

\subsection{Where it loads}

\reg{ES}:\reg{BX} gives the destination. With \reg{ES} = \hexa{0900} and \reg{BX} = 0, the physical
address is \hexa{9000}. Why there? Look back at the memory map in Chapter 4:

\begin{itemize}
  \item below \hexa{500} sit the interrupt vector table and BIOS data: untouchable;
  \item at \hexa{7C00} we are ourselves: we cannot write over our own running code;
  \item from \hexa{A0000} onwards begins video memory.
\end{itemize}

\hexa{9000} is comfortably free, leaves plenty of room below and above, and allows loading up to
\hexa{A0000} (some 60 KB) without hitting anything.

\begin{concept}{Why 60 sectors}
$60 \times 512 = 30{,}720$ bytes. The figure is hand-picked: enough for the current kernel
($\sim$23 KB) with headroom, and short of the 63-sectors-per-track limit that a single CHS read
imposes. If the kernel grew further, this number would have to grow \emph{and} you would have to
watch that the read does not cross a track boundary.
\end{concept}

\section{The final jump}

\begin{lstlisting}[style=asm]
jmp 0900h:0000
\end{lstlisting}

This is a \term{far jump}: besides changing \reg{IP}, it changes \reg{CS}. Changing both matters,
because the kernel code will compute its own addresses assuming \reg{CS} holds \hexa{0900}. From
this instruction on, the bootloader has finished its job and the kernel is in control.

\begin{recipe}{editing the boot sector}{ch12-bootloader}
The real bootloader with four marked edits: your own message, two characters printed with the BIOS
teletype call directly, and a pause using \cod{INT 15h / AH=86h} so you can read the result before
the kernel clears the screen.

It still loads the kernel and jumps to it, so the system boots normally afterwards. Watch the
\cod{times 510-(\$-\$\$)} line at the bottom: add too much and NASM fails with a negative count,
which is the 512-byte limit enforcing itself.
\end{recipe}

\begin{tryit}[Watch the real boot]
Build and run with \cod{make run}. You will see, in order:
\begin{lstlisting}[style=txt]
Welcome to the lytlnybl bootloader!
Loading up the kernel for you...
Hello World!, i am lytlnyblOS, in real mode
\end{lstlisting}
The first two are printed by the bootloader from \hexa{7C00}. The third is already printed by the
kernel from \hexa{9000}, still in real mode. If you see all three, the disk load worked.
\end{tryit}

% ============================================================
\chapter{The E820 memory map}

\begin{goals}
  \item Why the kernel cannot guess how much RAM exists
  \item How the map is requested from the BIOS
  \item How it is stored so the kernel can find it later
\end{goals}

\section{The problem}

The kernel needs to know which memory regions it may use. It cannot assume: every machine has a
different amount of RAM and holes in different places (memory reserved by the BIOS, by ACPI, by
memory-mapped devices).

There is only one moment when you can ask: \hi{while still in real mode}, because as soon as
we enter protected mode we lose access to BIOS services. That is why this job belongs to the
bootloader, even though the kernel is the one that needs the data.

\section{How you ask}

The service is called E820 and works by iteration: you call it in a loop and each call returns one
region, until it says it is done.

\filetag{lytlnybl/boot/bootloader.asm}
\begin{lstlisting}[style=asm]
store_memory_map:
    xor ax, ax
    mov ds, ax          ; DS = 0, to work with absolute addresses

    xor ebx, ebx        ; EBX = 0 means "first call"
    xor bp, bp          ; BP will count the entries

    mov di, memory_map_buffer   ; ES:DI = where to write (0x5000)
next_entry:
    mov ax, 0
    mov es, ax

    mov eax, 0xE820
    mov edx, 0x534D4150 ; 'SMAP' in ASCII: mandatory signature
    mov ecx, 24         ; size of each entry, in bytes

    int 15h
    jc done             ; carry = no more entries

    add di, 24          ; advance the write pointer
    inc bp              ; one more entry

    test ebx, ebx       ; did EBX return to 0? then that was the last
    jnz next_entry
done:
    mov [memory_map_entries], bp
    mov ax, 07C0h
    mov ds, ax          ; restore DS before returning
    ret
\end{lstlisting}

Details worth noticing:

\begin{description}
  \item[\reg{EBX} is the loop state] The BIOS uses it as a bookmark. You start at 0, it returns a
    value you must keep untouched for the next call, and when it returns to 0 the iteration is
    over. Touch \reg{EBX} between calls and you break the loop.
  \item[The \cod{'SMAP'} signature] \hexa{534D4150} is ``SMAP'' in ASCII. The BIOS checks it to
    distinguish this request from other \hexa{15} service functions.
  \item[\reg{DS} temporarily becomes 0] Because we want to write to absolute address \hexa{5000}.
    With \reg{DS}=0, \cod{di = 0x5000} points exactly there. It must be restored afterwards, or the
    rest of the bootloader prints garbage.
\end{description}

\section{The format of each entry}

Each entry occupies 24 bytes:

\begin{center}
\begin{tabularx}{0.95\textwidth}{@{}l l X@{}}
\toprule
\textbf{Offset} & \textbf{Size} & \textbf{Field} \\
\midrule
0  & 8 bytes & Base address of the region \\
8  & 8 bytes & Length in bytes \\
16 & 4 bytes & Type: \textbf{1 = usable}, anything else = reserved \\
20 & 4 bytes & Extended attributes (we do not use them) \\
\bottomrule
\end{tabularx}
\end{center}

And in C, the kernel reads it with a structure that mirrors that layout:

\filetag{lytlnybl/include/memory/pmm.h}
\begin{lstlisting}[style=c]
typedef struct {
    uint64_t base;
    uint64_t length;
    uint32_t type;
    uint32_t acpi_extended;
} __attribute__((packed)) memory_map_entry_t;
\end{lstlisting}

\begin{concept}{What \cod{\_\_attribute\_\_((packed))} does}
By default the compiler inserts padding between struct fields to align them and make access
faster. That is fine for internal structures, but disastrous when the structure must mirror a
layout defined by hardware or by another program. \cod{packed} tells the compiler: ``add no
padding, lay the fields out back to back''. It is used in every structure representing an external
format: GDT entries, IDT entries, the TSS, and the filesystem structures.
\end{concept}

\section{The agreement between bootloader and kernel}

The bootloader writes to two fixed addresses and the kernel reads them from there:

\begin{center}
\begin{tabularx}{0.9\textwidth}{@{}l l X@{}}
\toprule
\textbf{Address} & \textbf{Size} & \textbf{Contents} \\
\midrule
\hexa{4FFC} & 2 bytes & Number of entries in the map \\
\hexa{5000} & $24 \times n$ & The entries themselves \\
\bottomrule
\end{tabularx}
\end{center}

\filetag{lytlnybl/memory/pmm.c}
\begin{lstlisting}[style=c]
void init_pmm() {
    uint16_t map_entry_count = *(uint16_t*)0x4FFC;
    memory_map_entry_t* memory_map = (memory_map_entry_t*)0x5000;
    /* ... */
}
\end{lstlisting}

\begin{warn}{Coupling through magic numbers}
This is a fragile agreement: two different files, written in different languages, agreeing on a
constant typed by hand in both. If someone changes \hexa{5000} in the bootloader and forgets to
change it in \cod{pmm.c}, the kernel will read garbage as if it were a memory map and conclude that
it may use regions that do not exist.

Serious systems solve this with a documented boot protocol (Multiboot, for example) in which the
loader passes a pointer to a structure in a register. For a teaching system the shared constant is
acceptable --- but it is worth knowing it is technical debt.
\end{warn}

\begin{summary}
The bootloader has done its job: the kernel sits at \hexa{9000}, the memory map sits at \hexa{5000},
and control has just jumped into kernel code. But we are still in real mode, 16-bit, with no
protection at all. The next part changes that.
\end{summary}
